\chapter*{\fontsize{23.6pt}{28.3pt}\selectfont References}
\markboth{References}{References}
\addcontentsline{toc}{chapter}{References}

\begingroup
\setlength{\parindent}{0pt}
\raggedright

\par\hangindent=2em\hangafter=1 [Ama95] H. Amann. \emph{Gewöhnliche Differentialgleichungen}. W. de Gruyter, Berlin, 1983, 2. Aufl. 1995.\par

\par\hangindent=2em\hangafter=1 [Ape96] A. Apelblat. \emph{Tables of Integrals and Series}. Verlag Harri Deutsch, Frankfurt, 1996.\par

\par\hangindent=2em\hangafter=1 [Art31] E. Artin. \emph{Einführung in die Theorie der Gammafunktion}. Teubner, Leipzig, 1931.\par

\par\hangindent=2em\hangafter=1 [Art93] M. Artin. \emph{Algebra}. Birkhäuser, Basel, 1993.\par

\par\hangindent=2em\hangafter=1 [BBM86] A.P. Brudnikov, Yu.A. Brychkov, O.M. Marichev. \emph{Integrals and Series, I}. Gordon \& Breach, New York, 1986.\par

\par\hangindent=2em\hangafter=1 [BF87] M. Barner, F. Flohr. \emph{Analysis I, II}. W. de Gruyter, Berlin, 1987.\par

\par\hangindent=2em\hangafter=1 [Bla91] Ch. Blatter. \emph{Analysis I–III}. Springer Verlag, Berlin, 1991, 1992.\par

\par\hangindent=2em\hangafter=1 [Brü95] J. Brüdern. \emph{Einführung in die analytische Zahlentheorie}. Springer Verlag, Berlin, 1995.\par

\par\hangindent=2em\hangafter=1 [Car66] H. Cartan. \emph{Elementare Theorie der analytischen Funktionen einer oder mehrerer komplexen Veränderlichen}. BI Hochschultaschenbücher, 112/112a. Bibliographisches Institut, Mannheim, 1966.\par

\par\hangindent=2em\hangafter=1 [Con78] J.B. Conway. \emph{Functions of One Complex Variable}. Springer Verlag, Berlin, 1978.\par

\par\hangindent=2em\hangafter=1 [FB95] E. Freitag, R. Busam. \emph{Funktionentheorie}. Springer Verlag, Berlin, 1995.\par

\par\hangindent=2em\hangafter=1 [Gab96] P. Gabriel. \emph{Matrizen, Geometrie, Lineare Algebra}. Birkhäuser, Basel, 1996.\par

\par\hangindent=2em\hangafter=1 [GR81] I.S. Gradstein, I.M. Ryshik. \emph{Tables of Series, Products, and Integrals}. Verlag Harri Deutsch, Frankfurt, 1981.\par

\par\hangindent=2em\hangafter=1 [Koe83] M. Koecher. \emph{Lineare Algebra und analytische Geometrie}. Springer Verlag, Berlin, 1983.\par

\par\hangindent=2em\hangafter=1 [Kön92] K. Königsberger. \emph{Analysis 1, 2}. Springer Verlag, Berlin, 1992, 1993.\par

\par\hangindent=2em\hangafter=1 [Pra78] K. Prachar. \emph{Primzahlverteilung}. Springer Verlag, Berlin, 1978.\par

\par\hangindent=2em\hangafter=1 [Rem92] R. Remmert. \emph{Funktionentheorie 1, 2}. Springer Verlag, Berlin, 1992.\par

\par\hangindent=2em\hangafter=1 [Sch65] L. Schwartz. \emph{Méthodes Mathématiques pour les Sciences Physiques}. Hermann, Paris, 1965.\par

\par\hangindent=2em\hangafter=1 [Sch69] W. Schwarz. \emph{Einführung in die Methoden und Ergebnisse der Primzahltheorie}. BI Hochschultaschenbücher, 278/278a. Bibliographisches Institut, Mannheim, 1969.\par

\par\hangindent=2em\hangafter=1 [SS88] G. Scheja, U. Storch. \emph{Lehrbuch der Algebra}. Teubner, Stuttgart, 1988.\par

\par\hangindent=2em\hangafter=1 [Wal92] W. Walter. \emph{Analysis 1, 2}. Springer Verlag, Berlin, 1992.\par
\endgroup
